\section{Qué aprender del video a escala de Internet: representation, reward y action}

Después de MineCLIP, Voyager y Eureka, la clase pasa a direcciones de investigación de más largo plazo. El video de Internet es barato, dinámico y abarca una gran cantidad de conducta humana, pero proviene de perspectivas en tercera persona o no controladas y por lo general carece de etiquetas de agent action, state o grounding. El problema no debe llamarse de forma genérica «imitar a partir de video»; conviene descomponerlo en tres objetivos de supervisión: aprender una representation, aprender una reward, o inferir la action para luego hacer behavior cloning.

\subsection{Perspectiva: dos rutas hacia el foundation agent}

El ponente divide las direcciones futuras en dos tipos: seguir aprendiendo de internet-scale video y construir un multimodal foundation model nativo. La primera amplía los priors de conducta; la segunda unifica las interfaces de texto, imagen, video, sonido y action. Ambas deben volver, al final, al active loop y someterse a la retroalimentación de la embodiment y del environment.

\lecturefigure{slide-43-looking-forward-title.jpg}{Transición de los sistemas actuales a la dirección de largo plazo del generalist agent}{Diapositiva V043 recuperada de la grabación oficial; 00:37:36}

Esta página funciona como portada de sección: los sistemas anteriores ya demostraron que reward, code y memory son viables; lo que sigue es preguntar cómo ampliar los datos y las modalidades. Al leer la figura no hay que extrapolar directamente el éxito de chat/image/music a los agents: los datos de acción contienen causalidad, control y riesgo, y no se obtienen automáticamente con más tokens estáticos.

\lecturefigure{slide-44-two-foundation-directions.jpg}{Internet-scale video learning y multimodal foundation models}{Diapositiva V044 recuperada de la grabación oficial; 00:38:09}

\begin{importantbox}{Las dos rutas convergen en el action grounding}
El video learning aporta conducta dinámica y estrategias humanas; el multimodal model aporta una token/interface unificada. Aun así, el sistema necesita saber «qué action ejecuta quién y qué consecuencias tiene para la embodiment actual». Sin active grounding, el modelo aprende, como mucho, una observational correlation.
\end{importantbox}

\subsection{Ventajas y carencias del video de Internet}

El video está más cerca de la conducta que el texto, porque muestra cómo cambian los objetos con el tiempo, cómo la mano entra en contacto con las herramientas y cómo se desarrolla una tarea. Sin embargo, quien observa no sabe qué ocurre fuera de cuadro, ni ve el action command preciso, la fuerza, el tacto o los reintentos tras un fallo. La \term{behavior cloning (clonación de comportamiento)} entrena una policy con supervisión a partir de los pares observation--action de las demostraciones; cuando no hay action label, primero hay que recuperar o aproximar esos pares.

\lecturefigure{slide-45-internet-video-properties.jpg}{Internet video: barato, dinámico y rico, pero sin action ni grounding}{Diapositiva V045 recuperada de la grabación oficial; 00:38:48}

\begin{knowledgebox}{Términos clave: tres rutas para aprender del video}
\begin{tabularx}{\textwidth}{@{}L{0.19\textwidth}X X@{}}
\toprule
Ruta & Qué aprende & Carencia principal \\
\midrule
Representation & Representaciones espaciotemporales y semánticas transferibles & Una buena representación no implica saber controlar \\
Reward learning & Similitud/valor entre goal y trajectory & Reward bias y shortcut \\
Action / behavior cloning & Predecir una action ejecutable a partir de la observation & El video suele carecer de etiquetas de acción y la embodiment es distinta \\
\bottomrule
\end{tabularx}
\end{knowledgebox}

\lecturefigure{slide-46-internet-video-example.jpg}{El video de Internet en tercera persona muestra la conducta, pero oculta las señales de control}{Diapositiva V046 recuperada de la grabación oficial; 00:39:39}

En este ejemplo conviene mirar primero la secuencia temporal entre la mano, el objeto y el resultado, y después preguntar dónde está la action real. Los píxeles solo muestran el movimiento superficial; no proporcionan comandos articulares, fuerza aplicada, intención ni el estado fuera de cuadro, y una misma trayectoria visual puede ser producida por distintos controles. El video sirve para aprender semantic dynamics, pero requiere supuestos adicionales para convertirse en motor supervision de un robot concreto.

\teachervoice{El ponente describe el video de Internet como «cheap, dynamic, rich in human behavior» e inmediatamente añade que un observador pasivo no obtiene action labels ni grounding. Este giro es importante: la escala de los datos resuelve la coverage, pero no resuelve automáticamente el causal control; aquí reaparece el problema del gatito activo.}

\subsection{Ruta uno: Video Representation Learning}

La primera ruta aprende primero una representation espaciotemporal y después aplica fine-tune con una pequeña cantidad de agent data específica del dominio. Mediante future prediction, masked modeling o video--text alignment, el modelo puede obtener representaciones de objetos, fases de acción y cambios de escena, pero la capacidad de control sigue dependiendo de la action data y de la interface posteriores.

\lecturefigure{slide-47-representation-from-video.jpg}{Preentrenar la representation con video y luego adaptarla a tareas del agent}{Diapositiva V047 recuperada de la grabación oficial; 00:40:18}

\begin{knowledgebox}{Lectura de la figura: el preentrenamiento y el ajuste para control cumplen funciones distintas}
La etapa de video a gran escala de la figura aporta rasgos dinámicos generales; solo la pequeña cantidad posterior de agent data conecta la representation con un action space concreto. Primero compara el encoder congelado, el fine-tune y el entrenamiento desde cero, y luego revisa el tamaño de los datos posteriores; si no hay una action-grounded evaluation, solo se puede afirmar que la feature es útil, no que la policy haya aprendido la tarea.
\end{knowledgebox}

\subsection{Ruta dos: aprender la reward a partir del video}

La ruta anterior solo exige que el video preentrene una representación útil; esta sección va más allá y hace que el video participe directamente en la optimización de la tarea: MineCLIP convierte «cuánto se parece la trayectoria al objetivo» en una reward que guía al agent a explorar acciones en su propio entorno. No exige que el video y el robot compartan un control preciso, por lo que cruza interfaces con más facilidad que la imitación acción por acción; el costo es que hay que verificar si la semántica visual se transfiere y evitar que la policy explote activamente los puntos ciegos del reward model.

\lecturefigure{slide-48-reward-from-video.jpg}{Las representaciones video-language y goal-conditioned pueden formar una learned reward}{Diapositiva V048 recuperada de la grabación oficial; 00:40:45}

\begin{knowledgebox}{Lectura de la figura: la reward transfer es más débil, pero más flexible, que la action transfer}
El video solo necesita indicarle al sistema «si el resultado se ve más cerca del objetivo»; la action concreta la explora la policy en su propia embodiment, por lo que puede cruzar interfaces de control. Pero cuando la reward solo restringe la outcome appearance, puede pasar por alto la trayectoria, la seguridad o condiciones implícitas. Al leer los resultados hay que revisar a la vez el final state, la trajectory y la constraint.
\end{knowledgebox}

\subsection{Ruta tres: recuperar la action y hacer behavior cloning}

La tercera ruta es la más cercana al imitation learning: inferir una latent action a partir del video o usar un modelo para etiquetar la action, y luego entrenar una behavior policy. Video PreTraining (VPT), en Minecraft, entrena un inverse dynamics model con una pequeña cantidad de datos con etiquetas de teclado y ratón, y después completa la action de una gran cantidad de video sin etiquetar; este método depende de que la perspectiva de observación, el espacio de acciones y la embodiment coincidan lo suficiente.

\lecturefigure{slide-49-behavior-cloning-from-video.jpg}{Recuperar action labels a partir del video y luego hacer behavior cloning}{Diapositiva V049 recuperada de la grabación oficial; 00:42:33}

Si $z_t$ denota la latent action inferida a partir de cuadros adyacentes y la agent action real es $a_t$, puede escribirse:
\[
z_t\sim q_\psi(z\mid o_t,o_{t+1}),\qquad
a_t\sim p_\omega(a\mid z_t,e),\qquad
\mathcal{L}_{\mathrm{BC}}=-\sum_t\log\pi_\theta(a_t\mid o_{\le t},g).
\]
Aquí $q_\psi$ es el inverse/latent-action model, $p_\omega$ mapea la latent action a una acción ejecutable de la embodiment $e$, $\pi_\theta$ es la behavior policy y $g$ es el objetivo. Si la embodiment, la camera o el action support del video y del agent no coinciden, $p_\omega$ se convierte en la principal fuente de error.

\begin{warningbox}{Observation similarity no equivale a action identifiability}
A partir de solo dos imágenes por lo general no es posible recuperar de forma única la fuerza aplicada, la velocidad ni los contactos ocultos; el video en tercera persona además presenta movimiento de cámara y oclusiones. La action reconstruction debe reportar su uncertainty y permitir que la downstream policy corrija mediante la interacción con el entorno, en lugar de tratar las pseudoetiquetas como ground truth.
\end{warningbox}

\teachervoice{El ponente separa explícitamente representation, reward y action learning en tres rutas para evitar que «entrenar agents con video» se vuelva un lema vacío. La necesidad de action labels crece de una a otra: representation es la más débil, reward queda en medio y behavior cloning es la más fuerte; el sistema debe elegir la interfaz según la supervisión que pueda obtener.}

\subsection{Resumen de la sección}

El internet video ofrece datos de conducta humana baratos, dinámicos y de amplia cobertura, pero carece de action labels y grounding. El representation learning aprende rasgos, el reward learning aprende la similitud con el objetivo y la behavior cloning intenta recuperar las acciones; las tres rutas difieren en intensidad de supervisión, supuestos de transferencia y modos de falla. La escala puede ampliar los priors, pero la active environment feedback sigue siendo indispensable para verificar la causalidad de las acciones.
